% Luc Destombe --- licence CC BY-NC-SA 4.0
% https://creativecommons.org/licenses/by-nc-sa/4.0/deed.fr
% Fichier autonome : compiler avec pdflatex, deux fois (signets, nombre de pages).
\documentclass[a4paper,10pt]{article}
\makeatletter
\newif\iflycee@niveaux
\lycee@niveauxtrue
\newif\iflycee@palatino
\lycee@palatinotrue

\RequirePackage[utf8]{inputenc}
\RequirePackage[T1]{fontenc}
\RequirePackage[french]{babel}
\RequirePackage{etoolbox}

\newcommand{\ShorthandsOff}{\@ifundefined{shorthandoff}{}{\shorthandoff{;:!?}}}
\newcommand{\ShorthandsOn}{\@ifundefined{shorthandon}{}{\shorthandon{;:!?}}}
\ShorthandsOff
\AtBeginDocument{\ShorthandsOn}
\AtBeginEnvironment{tikzpicture}{\ShorthandsOff}

\RequirePackage{geometry}
\RequirePackage{fancyhdr}
\RequirePackage{lastpage}
\RequirePackage{multicol}
\RequirePackage{array}
\RequirePackage{enumitem}
\RequirePackage{xcolor}
\RequirePackage{needspace}

\RequirePackage{amsmath,amssymb}
\RequirePackage{mathpazo}
\DeclareMathSizes{10.95}{10.95}{8}{5.5}
\begingroup\catcode`\_=\active \catcode`\^=\active
\gdef_#1{\sb{\scriptscriptstyle #1}}
\gdef^#1{\sp{\scriptscriptstyle\lycee@moinsepais #1}}
\endgroup
\newcommand\lycee@moins{\mathbin{%
  \setbox\z@\hbox{$\m@th\scriptscriptstyle\lycee@moinsorig$}%
  \dimen@=.55\wd\z@ \advance\dimen@-.5pt
  \hbox to.75\wd\z@{\hss$\m@th\scriptscriptstyle\vcenter{\hbox{%
    \tikz\draw[line width=.5pt,line cap=round](0,0)--(\the\dimen@,0);}}$\hss}}}
\begingroup\lccode`\~=`\- \lowercase{\endgroup\let~}\lycee@moins
\newcommand\lycee@moinsepais{\mathcode`\-="8000 }
\AtBeginDocument{\mathchardef\lycee@moinsorig=\mathcode`\-\relax}
\AtBeginDocument{\catcode`\_=12 \mathcode`\_="8000
  \catcode`\^=12 \mathcode`\^="8000 }
\newcommand\lycee@exposants{%
  \fontdimen13\textfont2=.48\fontdimen6\textfont2
  \fontdimen14\textfont2=.48\fontdimen6\textfont2
  \fontdimen15\textfont2=.48\fontdimen6\textfont2
  \fontdimen13\scriptfont2=.48\fontdimen6\scriptfont2
  \fontdimen14\scriptfont2=.48\fontdimen6\scriptfont2
  \fontdimen15\scriptfont2=.48\fontdimen6\scriptfont2}
\every@math@size\expandafter{\the\every@math@size\lycee@exposants}
\newlength{\retraitcalculs}
\setlength{\retraitcalculs}{2em}

\RequirePackage{tikz}
\usetikzlibrary{arrows.meta,calc,babel,shapes.misc}
\RequirePackage[most]{tcolorbox}
\tcbuselibrary{skins,breakable}

\definecolor{coulDef}{HTML}{1F4E79}
\definecolor{coulProp}{HTML}{7030A0}
\definecolor{coulRem}{HTML}{C55A11}
\definecolor{coulEx}{HTML}{2E7D32}
\definecolor{coulExo}{HTML}{B03A2E}
\definecolor{coulPart}{HTML}{1F4E79}
\definecolor{coulMeth}{HTML}{1B6E4F}
\definecolor{coulCourbe}{HTML}{0B5ED7}
\definecolor{coulCourbeB}{HTML}{C00000}
\definecolor{coulCourbeC}{HTML}{2E7D32}
\definecolor{coulGrille}{HTML}{8C8C8C}
\definecolor{coulRep}{HTML}{1B6E4F}
\definecolor{coulBar}{HTML}{2E7D32}

\newcommand{\R}{\mathbb{R}}
\newcommand{\Rs}{\mathbb{R}^{*}}
\renewcommand{\le}{\leqslant}
\renewcommand{\ge}{\geqslant}
\newcommand{\vect}[1]{\overrightarrow{#1}}
\newcommand{\norme}[1]{\left\lVert #1 \right\rVert}
\newcommand{\intervalle}[2]{\left[#1\,;#2\right]}
\RequirePackage{cancel}
\renewcommand{\CancelColor}{\color{coulExo}}
\newcommand{\fois}[1]{\times{\color{coulExo}#1}}
\newcommand{\barre}[1]{\cancel{#1}}

\tikzset{grille/.style={coulGrille,line width=0.4pt}}
\newcommand{\quadrillage}[4]{\draw[grille,step=1] (#1,#2) grid (#3,#4);}
\tikzset{
  croix/.style={cross out,draw,line width=0.6pt,inner sep=0pt,minimum size=4pt},
  croixdroite/.style={croix,rotate=45,line width=0.8pt},
}
\newcommand{\pt}[3]{\path (#1) node[croix]{} node[#2,font=\small,inner sep=2.5pt]{$#3$};}

\newcommand{\lycee@repere}[4]{%
  \draw[grille,step=1] (#1,#2) grid (#3,#4);
  \draw[-{Stealth[length=2mm]},thick] (#1,0)--({#3+0.4},0) node[below left=-1pt]{$x$};
  \draw[-{Stealth[length=2mm]},thick] (0,#2)--(0,{#4+0.4}) node[below left=-1pt]{$y$};
  \node[below left,font=\scriptsize,inner sep=1.5pt] at (0,0) {$0$};}
\newcommand{\lycee@gradx}[1]{\foreach \k/\l in {#1}{%
  \draw (\k,0.07)--(\k,-0.07) node[below,font=\scriptsize,inner sep=1.5pt]{$\l$};}}
\newcommand{\lycee@grady}[1]{\foreach \k/\l in {#1}{%
  \draw (0.07,\k)--(-0.07,\k) node[left,font=\scriptsize,inner sep=1.5pt]{$\l$};}}
\newcommand{\lycee@intff}[2]{\left[#1\,;#2\right]}
\newcommand{\lycee@intfo}[2]{\left[#1\,;#2\right[}
\newcommand{\lycee@intof}[2]{\left]#1\,;#2\right]}
\newcommand{\lycee@intoo}[2]{\left]#1\,;#2\right[}
\newcommand{\lycee@ens}[1]{\left\{#1\right\}}
\AtBeginDocument{%
  \providecommand{\repere}{\lycee@repere}%
  \providecommand{\gradx}{\lycee@gradx}%
  \providecommand{\grady}{\lycee@grady}%
  \providecommand{\Cf}{\mathcal{C}\sb{\scriptscriptstyle f}}%
  \providecommand{\Cg}{\mathcal{C}\sb{\scriptscriptstyle g}}%
  \providecommand{\intff}{\lycee@intff}%
  \providecommand{\intfo}{\lycee@intfo}%
  \providecommand{\intof}{\lycee@intof}%
  \providecommand{\intoo}{\lycee@intoo}%
  \providecommand{\ens}{\lycee@ens}}

\newlist{questions}{enumerate}{3}
\setlist[questions,1]{label=\arabic*\textdegree),leftmargin=*,itemsep=2pt,topsep=3pt}
\setlist[questions,2]{label=\alph*),leftmargin=*,itemsep=1pt,topsep=2pt}
\setlist[questions,3]{label=\roman*),leftmargin=*,itemsep=1pt,topsep=1pt}
\setlist[itemize]{label=\textbullet,leftmargin=*,itemsep=2pt,topsep=3pt}

\newcommand{\besoinplace}[1]{\par\penalty\@M\begingroup
  \dimen@=#1\relax\dimen@ii=\pagegoal\advance\dimen@ii-\pagetotal
  \ifdim\dimen@>\dimen@ii\ifdim\dimen@ii>\z@\vfil\fi\break\fi\endgroup}

\newcommand{\lycee@pied}{\thepage/\pageref{LastPage}}
\newcommand{\entete}[2]{%
  \pagestyle{fancy}\fancyhf{}%
  \fancyhead[L]{\footnotesize\color{coulPart}\textbf{#1}}%
  \fancyhead[R]{\footnotesize\color{coulPart}\textbf{#2}}%
  \fancyfoot[C]{\footnotesize\lycee@pied}%
  \renewcommand{\headrulewidth}{0.4pt}%
  \renewcommand{\headrule}{\hbox to\headwidth{%
    \color{coulPart!40}\leaders\hrule height \headrulewidth\hfill}}}

\RequirePackage{ccicons}
\newcommand{\lycee@licence}{{\scriptsize\color{gray}%
  \copyright\ Luc Destombe --- \ccbyncsa\ CC BY-NC-SA 4.0}}
\newif\iflycee@licence \lycee@licencetrue
\newcommand{\sanslicence}{\lycee@licencefalse}
\AtBeginDocument{\iflycee@licence\AddToHookNext{shipout/foreground}{%
  \put(0,0){\rlap{\kern\dimexpr1in+\hoffset+\oddsidemargin+\textwidth\relax
    \llap{\raisebox{-\dimexpr1in+\voffset+\topmargin+\headheight+\headsep
      +\textheight+\footskip\relax}{\lycee@licence}}}}}\fi}

\newcommand{\formulesserrees}{%
  \setlength{\abovedisplayskip}{3pt}\setlength{\belowdisplayskip}{3pt}%
  \setlength{\abovedisplayshortskip}{2pt}\setlength{\belowdisplayshortskip}{3pt}}

\setlength{\parindent}{0pt}

\RequirePackage[implicit=false,hidelinks,pdfencoding=auto,psdextra,bookmarksopen,
  bookmarksopenlevel=1,pdfcreator={},pdfproducer={}]{hyperref}
\RequirePackage{bookmark}
\hypersetup{pdfauthor={Luc Destombe},pdfkeywords={CC BY-NC-SA 4.0}}
\newcounter{lycee@signet}
\newcommand{\signet}[3][]{\stepcounter{lycee@signet}%
  \edef\lycee@dest{\if\relax\detokenize{#1}\relax
    signet.\arabic{lycee@signet}\else#1\fi}%
  \hypertarget{\lycee@dest}{}\bookmark[level=#2,dest=\lycee@dest]{#3}}
\pdfstringdefDisableCommands{%
  \def\R{R}\def\vect#1{#1}\def\Cf{Cf}\def\Cg{Cg}\def\fois#1{×#1}%
  \def\barre#1{#1}\def\intervalle#1#2{[#1;#2]}\def\intff#1#2{[#1;#2]}%
  \def\textsuperscript#1{#1}\def\dfrac#1#2{#1/#2}\def\frac#1#2{#1/#2}%
  \def\vec#1{#1⃗}}%
\RequirePackage[official]{eurosym}

\geometry{top=0.8cm,bottom=1.3cm,left=1cm,right=1cm,footskip=0.6cm}

\pagestyle{fancy}
\fancyhf{}
\renewcommand{\headrulewidth}{0pt}
\fancyfoot[C]{\footnotesize Page \thepage{} / \pageref{LastPage}}

\newcommand{\titrefiche}[2]{%
  \begin{center}
    \Large\textbf{#1}\\[2pt]
    \large\textbf{#2}\\[2pt]
  \end{center}
  \vspace{0.10cm}}

\setlength{\columnseprule}{0.4pt}
\setlength{\columnsep}{15pt}
\renewcommand{\columnseprulecolor}{\color{black!45}}
\setlength{\multicolsep}{2pt}

\newcounter{partiefiche}
\renewcommand{\thepartiefiche}{\Roman{partiefiche}}
\newif\ifapresbandeau
\newcommand{\lycee@bandeau}[1]{%
  \par\penalty-600
  \vspace{0.08cm}%
  \noindent\colorbox{coulPart}{%
    \parbox{\dimexpr\linewidth-2\fboxsep\relax}{%
      \color{white}\bfseries\raggedright\strut#1\strut}}%
  \nopagebreak\par\nopagebreak
  \vspace{0.06cm}\nopagebreak
  \apresbandeautrue}
\newcommand{\partiefiche}[1]{%
  \refstepcounter{partiefiche}%
  \lycee@bandeau{\signet{1}{\thepartiefiche{} --- #1}\thepartiefiche{} --- #1}}
\newcommand{\partiefichesansnum}[1]{\lycee@bandeau{\signet{1}{#1}#1}}

\newcounter{exo}
\newlength{\espaceexo}\setlength{\espaceexo}{7pt}
\newlength{\espacetitre}\setlength{\espacetitre}{2pt}
\newcommand{\exo}[1][\relax]{%
  \par
  \ifapresbandeau\apresbandeaufalse\else\penalty-150\addvspace{\espaceexo}\fi
  \refstepcounter{exo}%
  \noindent\signet[exercice-\arabic{exo}]{2}{Exercice \arabic{exo}}%
  \textbf{\color{coulExo}\underline{Exercice \theexo}}%
  \ifx\relax#1\relax\else\textbf{ : #1}\fi
  \par\nopagebreak\vspace{\espacetitre}\nopagebreak
  \ignorespaces}

\setlist[enumerate]{nosep,leftmargin=*,itemsep=2pt,topsep=2pt}
\setlist[itemize]{nosep,leftmargin=*,topsep=2pt}
\newcommand{\choix}[3]{\par\hspace*{1em}%
  \makebox[0.3\linewidth][l]{a) #1}\makebox[0.3\linewidth][l]{b) #2}c) #3}
\newcommand{\deux}[2]{\par\hspace*{0.6em}%
  \makebox[0.48\linewidth][l]{#1}#2}
\newcommand{\trois}[3]{\par\hspace*{0.6em}%
  \makebox[0.32\linewidth][l]{#1}\makebox[0.32\linewidth][l]{#2}#3}

  \renewcommand{\exo}[1][\relax]{%
    \ifapresbandeau\apresbandeaufalse\else\par\penalty-150\fi
    \refstepcounter{exo}%
    \vspace{0.05cm}%
    \noindent\signet[exercice-\arabic{exo}]{2}{Exercice \arabic{exo}}%
    \textbf{\color{coulExo}\underline{Exercice \theexo}}%
    \ifx\relax#1\relax\else\textbf{ : #1}\fi%
    \par\nopagebreak
    \ignorespaces}
  \newcommand{\rappel}[1]{%
    \par\vspace{0.06cm}\nopagebreak
    \noindent\fcolorbox{black!35}{black!5}{%
      \parbox{\dimexpr\linewidth-2\fboxsep-2\fboxrule\relax}{%
        \small\textbf{Rappel.} #1}}%
    \par\vspace{0.06cm}\nopagebreak}
  \setlist[enumerate]{nosep,leftmargin=*}
  \setlist[itemize]{nosep,leftmargin=*}
  \setlength{\multicolsep}{3pt plus 1pt minus 1pt}
  \setlength{\premulticols}{10pt}
  \setlength{\postmulticols}{10pt}
  \newlist{expr}{enumerate}{1}
  \setlist[expr]{label={\Alph*\ $=$},nosep,leftmargin=*,itemsep=0pt}

\renewenvironment{center}{\par\addvspace{3pt}\centering}{\par\addvspace{3pt}}
\formulesserrees
\AtBeginDocument{\formulesserrees}
\clubpenalty=10000
\widowpenalty=10000
\displaywidowpenalty=10000
\brokenpenalty=10000
\setlength{\parskip}{0pt}

\RequirePackage{ragged2e}
\setlength{\RaggedRightRightskip}{0pt plus 1fil}
\AtBeginDocument{\RaggedRight\binoppenalty=10000 \relpenalty=10000 }
\g@addto@macro\@parboxrestore{\RaggedRight}
\makeatother
\usepackage{tkz-tab}
\makeatletter
\@namedef{nomfonc@1}{f}\@namedef{nomfonc@2}{g}\@namedef{nomfonc@3}{h}
\@namedef{nomfonc@4}{k}\@namedef{nomfonc@5}{l}\@namedef{nomfonc@6}{m}
\@namedef{nomfonc@7}{n}\@namedef{nomfonc@8}{p}\@namedef{nomfonc@9}{q}
\makeatother
\DeclareRobustCommand{\nomfonc}{\csname nomfonc@\number\value{expri}\endcsname}
\setlist[expr]{label={$\nomfonc(x)=$},nosep,leftmargin=*,widest*=8,itemsep=0pt}

\begin{document}

\titrefiche{1\textsuperscript{re} STI2D --- Polynômes du second degré}{Fiche d'exercices du chapitre}

\begin{multicols}{2}\raggedcolumns

\partiefiche{Fonction polynôme du second degré}

\exo[Est-ce un trinôme ?]
Pour chaque fonction, dire si c'est une fonction polynôme du second degré. Si oui,
donner $a$, $b$ et $c$ ; sinon, dire pourquoi.
\begin{multicols}{2}
\begin{enumerate}[label=\alph*)]
  \item $f(x)=3x^{2}-2x+5$
  \item $g(x)=4x+1$
  \item $h(x)=-x^{2}+6$
  \item $k(x)=2x^{3}+x^{2}-1$
  \item $l(x)=(x+2)(x-4)$
  \item $m(x)=x^{2}-(x-1)^{2}$
\end{enumerate}
\end{multicols}

\exo[Lire les coefficients]
Donner $a$, $b$ et $c$, \emph{avec leur signe}.
\begin{multicols}{2}
\begin{expr}
  \item $3x^{2}-4x+2$
  \item $-x^{2}+5x$
  \item $x^{2}-16$
  \item $-2x^{2}+7$
  \item $\dfrac{1}{3}x^{2}+x-4$
  \item $6-3x-x^{2}$
  \item $-x^{2}+2x-5$
  \item $5x^{2}$
\end{expr}
\end{multicols}

\exo[Développer d'abord]
Développer et réduire, puis donner $a$, $b$ et $c$ quand c'est un trinôme.
\emph{Deux de ces expressions n'en sont pas.}
\begin{multicols}{2}
\begin{expr}
  \item $(x+3)(x-4)$
  \item $(3x-2)(x+1)$
  \item $(x-2)^{2}+5$
  \item $2(x+3)^{2}-1$
  \item $(x+3)^{2}-(x-3)^{2}$
  \item $(3x+1)^{2}-9x^{2}$
\end{expr}
\end{multicols}

\exo[Calculer des images]
\begin{enumerate}[label=\arabic*)]
  \item $f(x)=3x^{2}-2x+1$ : calculer $f(0)$, $f(1)$, $f(-2)$ et
        $f\!\left(\dfrac{1}{3}\right)$.
  \item $g(x)=-x^{2}+6x$ : calculer $g(2)$, $g(-1)$, $g(7)$ et
        $g\!\left(\dfrac{1}{2}\right)$.
  \item Que remarque-t-on pour $g(-1)$ et $g(7)$ ?
\end{enumerate}

\partiefiche{Forme factorisée et racines}

\exo[De la forme factorisée à la forme développée]
Développer et réduire.
\begin{multicols}{2}
\begin{expr}
  \item $3(x-2)(x+1)$
  \item $-(x+3)(x-4)$
  \item $2(x+1)^{2}$
  \item $-3(x-1)(x-2)$
\end{expr}
\end{multicols}

\exo[Équations produit]
Résoudre dans $\R$ :
\begin{multicols}{2}
\begin{enumerate}[label=\alph*)]
  \item $(x-5)(x+2)=0$
  \item $-3(x+4)(2x-3)=0$
  \item $5(x+1)^{2}=0$
  \item $x(2x-8)=0$
\end{enumerate}
\end{multicols}

\exo[Lire les racines]
Donner les racines de chaque fonction, \emph{sans calcul}.
\begin{multicols}{2}
\begin{expr}
  \item $(x-3)(x+4)$
  \item $2(x+2)(x-6)$
  \item $-3(x-5)(x+1)$
  \item $(3x-1)(x+2)$
  \item $4(x-1)^{2}$
  \item $-(x+5)(x-5)$
  \item $x(x-7)$
  \item $2(x+1)(x+6)$
\end{expr}
\end{multicols}

\exo[Les deux formes d'une même fonction]
Soit $f$ la fonction définie sur $\R$ par $f(x)=3(x+1)(x-2)$.
\begin{enumerate}[label=\arabic*)]
  \item Développer et réduire $f(x)$.
  \item Calculer $f(0)$ avec chacune des deux formes.
  \item Donner les racines de $f$, puis résoudre $f(x)=0$.
  \item $f$ est-elle une fonction polynôme du second degré ? Donner $a$.
\end{enumerate}

\partiefiche{Le discriminant}

\exo[Calculer un discriminant]
Calculer $\Delta$ et donner son signe.
\begin{multicols}{2}
\begin{expr}
  \item $x^{2}-2x-15$
  \item $x^{2}+2x-15$
  \item $x^{2}+2x+15$
  \item $-x^{2}+2x-15$
  \item $4x^{2}-12x+9$
  \item $25x^{2}-10x+1$
  \item $-2x^{2}+5x-3$
  \item $-3x^{2}+4x-2$
  \item $-2x^{2}+7x-3$
\end{expr}
\end{multicols}

\exo[Combien de racines ?]
Calculer $\Delta$, puis dire combien le polynôme a de racines.
\begin{multicols}{2}
\begin{expr}
  \item $x^{2}+5x+4$
  \item $x^{2}+4x+5$
  \item $x^{2}-8x+7$
  \item $9x^{2}+6x+1$
  \item $2x^{2}+5x+3$
  \item $-3x^{2}-x-2$
\end{expr}
\end{multicols}

\exo[Quand $b$ ou $c$ est nul]
Même consigne. \emph{Attention aux termes absents et aux signes.}
\begin{multicols}{2}
\begin{expr}
  \item $x^{2}-25$
  \item $-3x^{2}+12$
  \item $2x^{2}-7x$
  \item $-x^{2}+6x$
  \item $-2x^{2}-5$
  \item $4x^{2}+1$
\end{expr}
\end{multicols}

\partiefiche{Racines et équations du second degré}

\exo[Deux racines]
Résoudre dans $\R$. \emph{Le discriminant est à chaque fois un carré parfait.}
\begin{multicols}{2}
\begin{enumerate}[label=\alph*)]
  \item $x^{2}-2x-15=0$
  \item $2x^{2}+3x-2=0$
  \item $3x^{2}-8x+5=0$
  \item $-x^{2}+7x-12=0$
\end{enumerate}
\end{multicols}

\exo[Les trois cas]
Résoudre dans $\R$, en précisant à chaque fois le cas.
\begin{multicols}{2}
\begin{enumerate}[label=\alph*)]
  \item $4x^{2}-20x+25=0$
  \item $x^{2}+2x+4=0$
  \item $-3x^{2}+2x-5=0$
  \item $x^{2}+8x+16=0$
  \item $6x^{2}+x-2=0$
  \item $-2x^{2}+5x-2=0$
\end{enumerate}
\end{multicols}

\exo[Sans discriminant]
Résoudre dans $\R$ \emph{sans calculer de discriminant}.
\begin{multicols}{2}
\begin{enumerate}[label=\alph*)]
  \item $x^{2}-7=0$
  \item $2x^{2}-9x=0$
  \item $x^{2}-64=0$
  \item $9x^{2}-4=0$
  \item $3x^{2}+12x=0$
  \item $x^{2}+9=0$
\end{enumerate}
\end{multicols}

\exo[Tout ramener d'un côté]
Résoudre dans $\R$ :
\begin{multicols}{2}
\begin{enumerate}[label=\alph*)]
  \item $3x^{2}+2x-1=2x^{2}+7$
  \item $x(x+3)=-(2x+6)$
  \item $x^{2}-3x+7=5$
  \item $(x+2)(x-3)=6$
\end{enumerate}
\end{multicols}

\exo[Racines d'une fonction]
Déterminer les racines de chaque fonction, si elle en a.
\begin{multicols}{2}
\begin{enumerate}[label=\alph*)]
  \item $f(x)=x^{2}-9x+20$
  \item $g(x)=-3x^{2}+3x+18$
  \item $h(x)=x^{2}+4x+7$
  \item $k(x)=9x^{2}-6x+1$
\end{enumerate}
\end{multicols}

\partiefiche{Factoriser un trinôme}

\exo[Les racines sont données]
Écrire la forme factorisée. \emph{Ne pas oublier le coefficient $a$.}
\begin{enumerate}[label=\alph*)]
  \item $x^{2}+x-6$, de racines $-3$ et $2$
  \item $3x^{2}+3x-18$, de racines $-3$ et $2$
  \item $-x^{2}+5x-4$, de racines $1$ et $4$
  \item $2x^{2}-12x+18$, de racine double $3$
\end{enumerate}

\exo[Factoriser]
Calculer $\Delta$, puis factoriser.
\begin{multicols}{2}
\begin{expr}
  \item $x^{2}-2x-8$
  \item $2x^{2}-x-1$
  \item $-x^{2}-2x+3$
  \item $3x^{2}-7x+2$
\end{expr}
\end{multicols}

\exo[Factoriser si c'est possible]
Factoriser quand c'est possible ; sinon, justifier.
\begin{multicols}{2}
\begin{expr}
  \item $x^{2}+4x+4$
  \item $x^{2}+9$
  \item $x^{2}-9$
  \item $6x^{2}+x-2$
  \item $3x^{2}+x+2$
  \item $-3x^{2}+6x-3$
\end{expr}
\end{multicols}

\exo[Racine évidente]
Pour chaque trinôme $f$, $g$, $h$, vérifier que le nombre indiqué est une racine, puis
factoriser.
\begin{enumerate}[label=\alph*)]
  \item $f(x)=x^{2}-6x+5$, avec $1$
  \item $g(x)=2x^{2}+5x-7$, avec $1$
  \item $h(x)=-2x^{2}-5x-3$, avec $-1$
\end{enumerate}

\partiefiche{Représentation graphique}

\exo[Tableau de valeurs et tracé]
Soit $f$ la fonction définie sur $\intff{-3}{2}$ par $f(x)=x^{2}+x-2$.
\begin{enumerate}[label=\arabic*)]
  \item Quel est le signe de $a$ ? La parabole est-elle un creux ou une bosse ?
  \item Calculer les racines de $f$.
  \item Recopier et compléter le tableau de valeurs :
\end{enumerate}
\begin{center}
\renewcommand{\arraystretch}{1.25}
\begin{tabular}{|c|c|c|c|c|c|c|}
\hline
$x$ & $-3$ & $-2$ & $-1$ & $0$ & $1$ & $2$\\
\hline
$f(x)$ & & & & & & \\
\hline
\end{tabular}
\end{center}
\begin{enumerate}[label=\arabic*),start=4]
  \item Placer les points et tracer la courbe dans le repère ci-dessous.
\end{enumerate}
\begin{center}
\begin{tikzpicture}[x=0.6cm,y=0.6cm]
  \repere{-4}{-3}{3}{5}
  \gradx{-3,-2,-1,1,2} \grady{-2,-1,1,2,3,4}
\end{tikzpicture}
\end{center}

\exo[Un deuxième tracé]
Même travail avec la fonction $g$ définie sur $\intff{0}{4}$ par $g(x)=-x^{2}+4x-3$ :
signe de $a$, racines, tableau de valeurs \emph{par pas de $1$}, puis tracé dans un
repère orthonormé.

\exo[Associer une courbe à une fonction]
On donne quatre fonctions :
\[
  f(x)=x^{2}-1 \qquad g(x)=-x^{2}+3
\]
\[
  h(x)=x^{2}+1 \qquad k(x)=(x+1)^{2}
\]
\begin{enumerate}[label=\arabic*)]
  \item Pour chacune, donner le signe de $a$ et le nombre de racines.
  \item Associer chaque fonction à sa courbe.
\end{enumerate}
\begin{center}
\begin{tikzpicture}[x=0.28cm,y=0.28cm]
  \repere{-3}{-5}{3}{5}
  \draw[coulCourbe,line width=1pt,domain=-2.8:2.8,samples=40] plot (\x,{-\x*\x+3});
  \node[below,font=\footnotesize] at (0,-5) {\textbf{Courbe 1}};
\end{tikzpicture}
\hspace{0.1cm}
\begin{tikzpicture}[x=0.28cm,y=0.28cm]
  \repere{-3}{-5}{3}{5}
  \draw[coulCourbe,line width=1pt,domain=-3:1.2,samples=40] plot (\x,{(\x+1)*(\x+1)});
  \node[below,font=\footnotesize] at (0,-5) {\textbf{Courbe 2}};
\end{tikzpicture}
\hspace{0.1cm}
\begin{tikzpicture}[x=0.28cm,y=0.28cm]
  \repere{-3}{-5}{3}{5}
  \draw[coulCourbe,line width=1pt,domain=-2.4:2.4,samples=40] plot (\x,{\x*\x-1});
  \node[below,font=\footnotesize] at (0,-5) {\textbf{Courbe 3}};
\end{tikzpicture}
\hspace{0.1cm}
\begin{tikzpicture}[x=0.28cm,y=0.28cm]
  \repere{-3}{-5}{3}{5}
  \draw[coulCourbe,line width=1pt,domain=-2:2,samples=40] plot (\x,{\x*\x+1});
  \node[below,font=\footnotesize] at (0,-5) {\textbf{Courbe 4}};
\end{tikzpicture}
\end{center}

\exo[Sans tracer]
Dire, en une ligne de calcul, en combien de points la parabole coupe l'axe des
abscisses.
\begin{multicols}{2}
\begin{expr}
  \item $x^{2}-8x+16$
  \item $3x^{2}+x+4$
  \item $-x^{2}+2x+8$
  \item $2x^{2}-18$
  \item $-9x^{2}+6x-1$
  \item $x^{2}+4x+6$
\end{expr}
\end{multicols}

\partiefiche{Sommet et variations}

\exo[Coordonnées du sommet]
Calculer $\alpha$ et $\beta$, puis donner les coordonnées du sommet $S$.
\begin{multicols}{2}
\begin{enumerate}[label=\alph*)]
  \item $f(x)=x^{2}-6x+2$
  \item $g(x)=-x^{2}+4x+3$
  \item $h(x)=2x^{2}+8x+5$
  \item $k(x)=-2x^{2}+12x-7$
\end{enumerate}
\end{multicols}

\exo[Tableau de variations]
Dresser le tableau de variations sur $\R$, puis dire si la fonction a un minimum ou un
maximum, et donner sa valeur.
\begin{multicols}{2}
\begin{enumerate}[label=\alph*)]
  \item $f(x)=x^{2}-4x+3$
  \item $g(x)=-3x^{2}+6x+9$
  \item $h(x)=2x^{2}+8x$
  \item $k(x)=-x^{2}-6x-5$
\end{enumerate}
\end{multicols}

\exo[Le sommet à partir des racines]
Soit $f(x)=x^{2}-6x+8$ et $g(x)=-2x^{2}+8x+10$.
\begin{enumerate}[label=\arabic*)]
  \item Déterminer les racines de $f$, puis celles de $g$.
  \item En déduire $\alpha$ pour chacune, sans la formule $-\dfrac{b}{2a}$.
        \emph{(Le sommet est à mi-chemin des racines.)}
  \item Calculer $\beta$, puis dresser les deux tableaux de variations.
\end{enumerate}

\exo[Maximum ou minimum ?]
Lesquelles de ces fonctions ont un maximum ? un minimum ? Donner sa valeur et le $x$ où
il est atteint.
\begin{multicols}{2}
\begin{enumerate}[label=\alph*)]
  \item $f_{1}(x)=4x^{2}-8x+1$
  \item $f_{2}(x)=-x^{2}+4x+2$
  \item $f_{3}(x)=-3x^{2}+6x$
  \item $f_{4}(x)=x^{2}+8x+12$
\end{enumerate}
\end{multicols}

\partiefiche{Synthèse : étudier et tracer une parabole}

\rappel{Étude complète : sens de la parabole (signe de $a$) ; sommet ; racines ;
tableau de variations ; tableau de valeurs ; tracé.}

\exo[Étude complète]
Faire l'étude complète de la fonction $f$ définie sur $\intff{-2}{4}$ par
$f(x)=x^{2}-2x-3$, puis tracer sa courbe dans un repère orthonormé.

\exo[Étude complète]
Même travail avec la fonction $g$ définie sur $\intff{-3}{2}$ par $g(x)=-x^{2}-x+2$.

\exo[Étude complète]
Même travail avec la fonction $h$ définie sur $\intff{-2}{1}$ par $h(x)=2x^{2}+2x-4$.
Vérifier que la courbe passe par les points d'ordonnée $0$ calculés.

\partiefiche{Signe d'un trinôme}

\exo[À partir de la forme factorisée]
Dresser le tableau de signes de chaque polynôme.
\begin{multicols}{2}
\begin{expr}
  \item $3(x+2)(x-5)$
  \item $-2(x-1)(x+6)$
  \item $4(x+4)(3x-1)$
  \item $-3(2x-5)(4x+1)$
\end{expr}
\end{multicols}

\exo[À partir de la forme développée]
Calculer $\Delta$, les racines, puis dresser le tableau de signes.
\begin{multicols}{2}
\begin{enumerate}[label=\alph*)]
  \item $f(x)=2x^{2}+5x-12$
  \item $g(x)=-x^{2}+2x+8$
  \item $h(x)=2x^{2}-7x+3$
  \item $k(x)=-3x^{2}-x+2$
\end{enumerate}
\end{multicols}

\exo[Les trois cas]
Étudier le signe de chaque polynôme.
\begin{multicols}{2}
\begin{expr}
  \item $x^{2}+4x+4$
  \item $x^{2}-4$
  \item $x^{2}+4$
  \item $-x^{2}+2x-4$
  \item $-9x^{2}+6x-1$
  \item $2x^{2}-5x+3$
\end{expr}
\end{multicols}

\exo[Sans tableau de signes]
Pour chaque trinôme, dire s'il est \emph{toujours positif}, \emph{toujours négatif}, ou
s'il \emph{change de signe}, en s'appuyant sur le signe de $a$ et celui de $\Delta$.
\begin{multicols}{2}
\begin{expr}
  \item $x^{2}+2x+3$
  \item $-2x^{2}-3$
  \item $3x^{2}-12$
  \item $-x^{2}+4x-4$
  \item $4x^{2}+x+1$
  \item $-3x^{2}+2x-1$
\end{expr}
\end{multicols}

\partiefiche{Inéquations du second degré}

\exo[Inéquations]
Résoudre dans $\R$, en donnant les solutions sous forme d'intervalles.
\begin{multicols}{2}
\begin{enumerate}[label=\alph*)]
  \item $x^{2}+4x+4>0$
  \item $-2x^{2}+5x-3\le 0$
  \item $x^{2}-2x-15\ge 0$
  \item $x^{2}+3x-10<0$
\end{enumerate}
\end{multicols}

\exo[Tout ramener d'un côté]
Résoudre dans $\R$ :
\begin{multicols}{2}
\begin{enumerate}[label=\alph*)]
  \item $-2x^{2}+7x\le 3$
  \item $2x^{2}\ge 3x-4$
  \item $x(x-4)\ge x-6$
  \item $x^{2}+8<6x$
\end{enumerate}
\end{multicols}

\exo[Cas particuliers]
Résoudre dans $\R$ :
\begin{multicols}{3}
\begin{enumerate}[label=\alph*)]
  \item $x^{2}\ge 25$
  \item $x^{2}<9$
  \item $-x^{2}+1>0$
  \item $3x^{2}+6\le 0$
  \item $(x+2)^{2}>0$
  \item $x^{2}-x+4>0$
\end{enumerate}
\end{multicols}

\exo[Lecture graphique]
Voici la courbe de la fonction $f$ définie par $f(x)=x^{2}+2x-3$.
\begin{center}
\begin{tikzpicture}[x=0.36cm,y=0.36cm]
  \repere{-5}{-5}{3}{5}
  \gradx{-4,-3,-2,-1,1,2} \grady{-4,-3,-2,-1,1,2,3,4}
  \draw[coulCourbe,line width=1.1pt,domain=-4:2,samples=60] plot (\x,{\x*\x+2*\x-3});
  \node[coulCourbe,font=\footnotesize] at (2.5,3.6) {$\Cf$};
\end{tikzpicture}
\end{center}
Résoudre \emph{graphiquement}, puis \emph{par le calcul} :
\begin{multicols}{2}
\begin{enumerate}[label=\alph*)]
  \item $f(x)=0$
  \item $f(x)>0$
  \item $f(x)\le 0$
  \item $f(x)=-3$
\end{enumerate}
\end{multicols}

\partiefiche{Intersection d'une parabole et d'une droite}

\exo[Points d'intersection]
Trouver les points d'intersection de $\mathcal{P}$ et $\mathcal{D}$, s'il y en a.
\begin{enumerate}[label=\alph*)]
  \item $\mathcal{P}$ : $y=x^{2}-4x+2$ \ et \ $\mathcal{D}$ : $y=-x+2$
  \item $\mathcal{P}$ : $y=x^{2}+2x+4$ \ et \ $\mathcal{D}$ : $y=4x+3$
  \item $\mathcal{P}$ : $y=x^{2}-x+3$ \ et \ $\mathcal{D}$ : $y=x+1$
\end{enumerate}

\exo[Deux paraboles]
Trouver les points d'intersection de $\mathcal{P}$ et $\mathcal{P}'$, s'il y en a.
\begin{enumerate}[label=\alph*)]
  \item $\mathcal{P}$ : $y=x^{2}+x-3$ \ et \ $\mathcal{P}'$ : $y=-x^{2}+3x+1$
  \item $\mathcal{P}$ : $y=x^{2}-2x+4$ \ et \ $\mathcal{P}'$ : $y=-x^{2}+2x+2$
  \item $\mathcal{P}$ : $y=3x^{2}-2x+1$ \ et \ $\mathcal{P}'$ : $y=3x^{2}+4x-5$
        \emph{(que remarque-t-on ?)}
\end{enumerate}

\exo[Combien de points ?]
$\mathcal{P}$ est la parabole d'équation $y=x^{2}$. Pour chaque droite, dire
\emph{sans calculer les points} combien elle a de points communs avec $\mathcal{P}$ :
\[
  \mathcal{D}_{1} : y=4x-4 \qquad \mathcal{D}_{2} : y=4x-5 \qquad \mathcal{D}_{3} : y=4x+1
\]

\exo[Avec un paramètre]
$\mathcal{P}$ est la parabole d'équation $y=x^{2}+x$ et $\mathcal{D}_{m}$ la droite
d'équation $y=3x+m$, où $m$ est un nombre réel.
\begin{enumerate}[label=\arabic*)]
  \item Montrer que les abscisses des points communs sont les solutions de
        $x^{2}-2x-m=0$.
  \item Calculer $\Delta$ en fonction de $m$.
  \item Pour quelles valeurs de $m$ la droite coupe-t-elle $\mathcal{P}$ en deux points ?
        en un seul ? en aucun ?
\end{enumerate}

\partiefiche{Position relative de deux courbes}

\exo[Une parabole et une droite]
Soit $f(x)=x^{2}+x-2$ et $g(x)=3x-2$.
\begin{enumerate}[label=\arabic*)]
  \item Calculer et réduire $d(x)=f(x)-g(x)$.
  \item Étudier le signe de $d(x)$.
  \item En déduire la position relative de $\Cf$ et de la droite $\mathcal{C}_{g}$.
\end{enumerate}

\exo[Trois situations]
Étudier la position relative de $\Cf$ et $\mathcal{C}_{g}$.
\begin{enumerate}[label=\alph*)]
  \item $f(x)=x^{2}-4x+5$ \ et \ $g(x)=x+1$
  \item $f(x)=-x^{2}+3x+1$ \ et \ $g(x)=x+1$
  \item $f(x)=x^{2}+3x+1$ \ et \ $g(x)=5x$
\end{enumerate}

\exo[Deux paraboles]
Étudier la position relative de $\Cf$ et $\mathcal{C}_{g}$.
\begin{enumerate}[label=\alph*)]
  \item $f(x)=x^{2}+x+1$ \ et \ $g(x)=-x^{2}+4x+3$
  \item $f(x)=3x^{2}+x-2$ \ et \ $g(x)=3x^{2}-5x+4$
  \item $f(x)=2x^{2}+3x-1$ \ et \ $g(x)=x^{2}+5x+2$
\end{enumerate}

\partiefiche{Problèmes du second degré}

\exo[Une aire imposée]
Un rectangle a pour périmètre $14$~cm. On note $x$ sa longueur, en cm.
\begin{enumerate}[label=\arabic*)]
  \item Entre quelles valeurs $x$ varie-t-il ?
  \item Exprimer la largeur, puis l'aire, en fonction de $x$.
  \item On veut une aire supérieure à $10~\text{cm}^{2}$. Écrire l'inéquation, la
        résoudre, et conclure.
\end{enumerate}

\exo[Une trajectoire]
Un projectile est lancé depuis le toit d'un bâtiment. Sa hauteur, en mètres, au bout de
$t$ secondes est
\[
  h(t)=-5t^{2}+30t+35 \qquad (t\ge 0).
\]
\begin{enumerate}[label=\arabic*)]
  \item De quelle hauteur part le projectile ?
  \item Quelle hauteur maximale atteint-il, et à quel instant ?
  \item À quel instant touche-t-il le sol ?
  \item Pendant combien de temps est-il plus haut que son point de départ ?
\end{enumerate}

\exo[Un atelier]
Un atelier fabrique $x$ pièces par jour, avec $0\le x\le 50$. Le coût de production, en
euros, est $C(x)=x^{2}-8x+400$, et la recette $R(x)=42x$.
\begin{enumerate}[label=\arabic*)]
  \item Montrer que le bénéfice est $B(x)=-x^{2}+50x-400$.
  \item Pour quelles productions l'atelier fait-il un bénéfice ?
  \item Quelle production rend le bénéfice maximal ? Combien vaut-il ?
\end{enumerate}

\end{multicols}

\end{document}
